% =========================================================================
%  RansomShield — Présentation de soutenance  (design orange moderne)
%  IFRI-UAC | Codjo Fréjus Loïc SANGRONIO | Ing. Guy KPEDJO | 2023-2024
%  Compilateur : pdfLaTeX
%
%  NOTES ORATEUR — DEUX PDF GÉNÉRÉS AUTOMATIQUEMENT, RIEN À DÉCOMMENTER :
%    ./build.sh
%      → presentation_soutenance.pdf        (PROPRE : à projeter — Meet, vidéoprojecteur)
%      → presentation_soutenance_notes.pdf  (AVEC NOTES : à garder sur votre propre écran)
%    Le choix se fait automatiquement selon le nom de compilation (\jobname) :
%    compiler ce fichier normalement ne montre JAMAIS les notes.
% =========================================================================

\documentclass[aspectratio=169,t,10pt]{beamer}

\usepackage[utf8]{inputenc}
\usepackage[T1]{fontenc}
\usepackage[french]{babel}
\usepackage{lmodern}
\usepackage{microtype}
\usepackage{xcolor}
\usepackage{graphicx}
\usepackage{tikz}
\usetikzlibrary{arrows.meta,positioning,fit,backgrounds,shadows.blur,calc,shapes.geometric}
\usepackage{fontawesome5}
\usepackage{booktabs}
\usepackage{tabularx}
\usepackage{array}
\usepackage{colortbl}
\usepackage{tcolorbox}
\tcbuselibrary{skins,breakable}
\usepackage{ragged2e}
\usepackage{adjustbox}
\usepackage{pgfpages}   % nécessaire pour les notes sur 2e écran

% ─── NOTES ORATEUR — activées SEULEMENT pour le job "..._notes" ──────────
%  Voir build.sh : deux compilations du même fichier, deux PDF différents.
\edef\rs@jobname{\jobname}
\ifnum\pdfstrcmp{\rs@jobname}{presentation_soutenance_notes}=0
  \setbeameroption{show notes on second screen=right}
\fi
% ─────────────────────────────────────────────────────────────────────────

% =========================================================================
%  PALETTE BLEU MARINE IFRI
% =========================================================================
\definecolor{primary}   {RGB}{  0,  82, 165}   % bleu IFRI principal
\definecolor{darkblue}  {RGB}{ 10,  40,  90}   % bleu nuit (ombres, titre slide final)
\definecolor{darkbg}    {RGB}{ 12,  28,  68}   % fond sombre (slide final)
\definecolor{light}     {RGB}{232, 240, 252}   % fond très clair bleu des blocs
\definecolor{muted}     {RGB}{100, 100, 110}   % gris neutre
\definecolor{crit}      {RGB}{185,  28,  28}   % rouge critique
\definecolor{high}      {RGB}{194,  96,  10}   % orange élevé
\definecolor{norm}      {RGB}{ 22, 128,  65}   % vert normal
\definecolor{textdark}  {RGB}{ 18,  22,  55}   % texte principal bleu sombre

% =========================================================================
%  THÈME BEAMER
% =========================================================================
\usetheme{default}
\useinnertheme{circles}

\setbeamercolor{structure}           {fg=primary}
\setbeamercolor{normal text}         {fg=textdark,bg=white}
\setbeamercolor{alerted text}        {fg=crit}
\setbeamercolor{frametitle}          {fg=primary,bg=white}
\setbeamerfont{frametitle}           {size=\large,series=\bfseries}
\setbeamercolor{block title}         {fg=white,bg=primary}
\setbeamercolor{block body}          {fg=textdark,bg=light}
\setbeamercolor{block title alerted} {fg=white,bg=crit}
\setbeamercolor{block body alerted}  {fg=textdark,bg=crit!8}
\setbeamerfont{block title}          {size=\small,series=\bfseries}
\setbeamerfont{block body}           {size=\small}
\setbeamercolor{item}                {fg=primary}
\setbeamertemplate{itemize item}     {\textcolor{primary}{\small\textbullet}}
\setbeamertemplate{itemize subitem}  {\textcolor{primary!70}{\footnotesize$\triangleright$}}
\setbeamertemplate{navigation symbols}{}

% Filet orange sous chaque titre de frame
\addtobeamertemplate{frametitle}{}{%
  \vspace{-4pt}%
  {\color{primary}\rule{\textwidth}{2pt}}%
  \vspace{2pt}%
}

% Filigrane bouclier (très discret, fond de chaque slide)
\setbeamertemplate{background}{%
  \begin{tikzpicture}[overlay,remember picture]
    \node[anchor=south east, opacity=0.045,
          inner sep=0pt]
      at ([xshift=-0.5cm,yshift=0.8cm]current page.south east)
      {\fontsize{90}{90}\selectfont\color{primary}\faIcon{shield-alt}};
  \end{tikzpicture}%
}

% Footline orange 3 colonnes
\setbeamertemplate{footline}{%
  \leavevmode%
  \begin{beamercolorbox}[wd=\paperwidth,ht=2.5ex,dp=1.2ex,colsep=0pt]{section in head/foot}%
    \color{white}%
    \hspace*{6pt}\tiny Codjo Fréjus Loïc SANGRONIO%
    \hfill%
    \tiny RansomShield~--- Détection \& Réponse ransomware%
    \hfill%
    \tiny 24 juin 2026\enspace\insertframenumber/\inserttotalframenumber\hspace*{6pt}%
  \end{beamercolorbox}%
}
\setbeamercolor{section in head/foot}{fg=white,bg=primary}

% =========================================================================
%  CHEMINS DES IMAGES
% =========================================================================
\graphicspath{%
  {New_IFRI_template_latex_05_04_2020_Bachelor___Licence__2_-2/images/}%
  {New_IFRI_template_latex_05_04_2020_Bachelor___Licence__2_-2/images/captures/}%
  {New_IFRI_template_latex_05_04_2020_Bachelor___Licence__2_-2/images/figures/}%
}

% =========================================================================
%  MACROS
% =========================================================================
\newcommand{\ok}  {\textcolor{norm}{\faCheckCircle}\ }
\newcommand{\bad} {\textcolor{crit}{\faTimesCircle}\ }
\newcommand{\warn}{\textcolor{high}{\faExclamationCircle}\ }
\newcommand{\arr} {\textcolor{primary}{\small\faChevronRight}\ }

% Slide séparateur de section (bande orange + numéro + icône thématique)
\newcommand{\sectionslide}[3]{%
  \begin{frame}[plain,noframenumbering]%
    \begin{tikzpicture}[overlay,remember picture]%
      % Bande orange côté gauche
      \fill[primary]
        (current page.north west) --
        ([xshift=5.5cm]current page.north west) --
        ([xshift=3.8cm]current page.south west) --
        (current page.south west) -- cycle;
      % Icône de section en filigrane (grand format, dans la bande)
      \node[anchor=center,
            font=\fontsize{95}{95}\selectfont,
            text=white,opacity=0.16]
        at ([xshift=2.0cm,yshift=1.7cm]current page.west) {\faIcon{#3}};
      % Numéro de section en filigrane
      \node[anchor=center,rotate=90,
            font=\fontsize{60}{60}\selectfont\bfseries,
            text=white,opacity=0.22]
        at ([xshift=2.0cm,yshift=-1.9cm]current page.west) {#1};
      % Icône + titre à droite
      \node[anchor=west, font=\LARGE\bfseries, text=textdark]
        at ([xshift=6.5cm,yshift=0.3cm]current page.west)
        {\textcolor{primary}{\faIcon{#3}}\ \ #2};
      % Trait orange sous le titre
      \draw[primary,line width=2.5pt]
        ([xshift=6.5cm,yshift=-0.4cm]current page.west) --
        ([xshift=14.5cm,yshift=-0.4cm]current page.west);
    \end{tikzpicture}%
  \end{frame}%
}

% =========================================================================
%  MÉTADONNÉES
% =========================================================================
\title[RansomShield]{%
  Conception et mise en place d'un système de\\[3pt]
  \textbf{détection et de réponse aux attaques ransomware}\\[3pt]
  dans un environnement contrôlé%
}
\author[L.~SANGRONIO]{Codjo Fréjus Loïc \textbf{SANGRONIO}}
\institute[IFRI-UAC]{%
  Licence Informatique --- Option Sécurité Informatique \textbar{} IFRI-UAC}
\date{24 juin 2026}

% =========================================================================
\begin{document}
% =========================================================================

% ─────────────────────────────────────────────────────────────────────────
%  SLIDE 1 — TITRE
% ─────────────────────────────────────────────────────────────────────────
{
\setbeamertemplate{headline}{}
\setbeamertemplate{footline}{}
\begin{frame}[plain,noframenumbering]

  \begin{tikzpicture}[overlay,remember picture]
    \node[anchor=north west,inner sep=8pt]
      at (current page.north west)
      {\includegraphics[height=1.35cm]{logoUAC1}};
    \node[anchor=north east,inner sep=8pt]
      at (current page.north east)
      {\includegraphics[height=1.35cm]{logoIfri}};
    % Footline titre
    \fill[primary]
      (current page.south west) rectangle
      ([yshift=1.9em]current page.south east);
    \node[anchor=west,font=\tiny\color{white},inner sep=6pt]
      at ([yshift=0.95em]current page.south west)
      {Codjo Fréjus Loïc SANGRONIO};
    \node[anchor=center,font=\tiny\color{white}]
      at ([yshift=0.95em]current page.south)
      {RansomShield --- Détection \& Réponse ransomware};
    \node[anchor=east,font=\tiny\color{white},inner sep=6pt]
      at ([yshift=0.95em]current page.south east)
      {24 juin 2026 \quad 1\,/\,\inserttotalframenumber};
  \end{tikzpicture}

  \vspace{1.4cm}

  \begin{tcolorbox}[
    colback=primary,colframe=darkblue,
    arc=10pt,boxrule=0pt,halign=center,
    left=16pt,right=16pt,top=12pt,bottom=12pt,
    drop shadow={opacity=0.3,shadow xshift=0pt,shadow yshift=-3pt}
  ]
    {\large\bfseries\color{white}
     Conception et mise en place d'un système de\\[6pt]
     détection et de réponse aux attaques ransomware\\[6pt]
     dans un environnement contrôlé}
  \end{tcolorbox}

  \vspace{0.45cm}

  \centering
  {\normalsize\bfseries Codjo Fréjus Loïc SANGRONIO}\\[4pt]
  {\small Licence Informatique --- Option Sécurité Informatique}\\[8pt]
  {\normalsize\bfseries 24 juin 2026}

  \vfill

  \begin{columns}[b]
    \column{0.5\textwidth}
      \raggedright{\scriptsize\color{muted}%
        \textbf{Encadreur :}\\Ing. Guy KPEDJO}
    \column{0.5\textwidth}
      \raggedleft{\scriptsize\color{muted}%
        \textbf{Maître de stage :}\\{[}Nom du maître de stage{]}}
  \end{columns}

  \vspace{0.45cm}
\end{frame}
}
\note{
  \textbf{[OUVERTURE — 0:00]}\\[4pt]
  Monsieur le Président du jury, Mesdames et Messieurs les membres du jury, bonjour.
  Je vous remercie pour l'occasion qui m'est donnée de présenter mon travail portant sur
  la conception et la mise en place d'un système de détection et de réponse aux attaques
  ransomware dans un environnement contrôlé~: RansomShield.
  Cette présentation durera dix minutes, suivie d'une démonstration live de cinq minutes
  sur des machines virtuelles réelles. Je reste à votre disposition pour vos questions
  à la fin.
}

% ─────────────────────────────────────────────────────────────────────────
%  SLIDE 2 — SOMMAIRE
% ─────────────────────────────────────────────────────────────────────────
\begin{frame}{Sommaire}
  \vspace{0.2cm}
  \begin{columns}[T]
    \column{0.58\textwidth}
      \begin{enumerate}\setlength\itemsep{10pt}
        \item[\textcolor{primary}{\bfseries\large 1.}] Introduction Générale
        \item[\textcolor{primary}{\bfseries\large 2.}] Revue de littérature
        \item[\textcolor{primary}{\bfseries\large 3.}] Analyse et Conception
        \item[\textcolor{primary}{\bfseries\large 4.}] Réalisation et Résultats
        \item[\textcolor{primary}{\bfseries\large 5.}] Conclusion et Perspectives
      \end{enumerate}
    \column{0.39\textwidth}
      \begin{tcolorbox}[
        colback=light,colframe=primary,arc=6pt,boxrule=1pt,
        title={\small\faClock\quad Informations pratiques},
        coltitle=white,fonttitle=\bfseries\small,
        attach boxed title to top left={yshift=-2mm,xshift=4mm},
        boxed title style={colback=primary,arc=3pt,boxrule=0pt}
      ]
        \small\color{textdark}
        \begin{tabular}{@{}r@{\;:\;}l@{}}
          Présentation & \textbf{10 min} \\[4pt]
          Diapositives & \textbf{21}     \\[4pt]
          Démo live    & \textbf{5 min}  \\[4pt]
          VMs test     & 3 KVM           \\
        \end{tabular}
      \end{tcolorbox}
  \end{columns}
\end{frame}
\note{
  \textbf{[SOMMAIRE — 0:30]}\\[4pt]
  Ma présentation s'articule en cinq parties~: le contexte et la problématique,
  la revue de littérature, la conception de la solution, sa réalisation avec les
  résultats obtenus, puis la conclusion. Je terminerai par une courte démonstration
  du système en fonctionnement.
}

% =========================================================================
\section{Introduction Générale}
% =========================================================================

\sectionslide{01}{Introduction Générale}{bullseye}

% ─── Contexte + Problématique ─────────────────────────────────────────────
\begin{frame}{Introduction Générale}

  \begin{block}{Contexte}
    L'évolution rapide du numérique expose les organisations aux ransomwares~: des
    logiciels malveillants qui rendent les données inaccessibles puis exigent une rançon.
    \begin{itemize}\setlength\itemsep{2pt}
      \item \textbf{+73\,\%} d'attaques signalées en 2023 \textit{(ENISA)}
      \item Délai moyen de détection \textbf{sans outil dédié~: 21 jours}
      \item Coût moyen par attaque~: \textbf{$>$ 1 million de dollars}
      \item Cibles~: PME, hôpitaux, administrations, universités
    \end{itemize}
  \end{block}

  \vspace{4pt}

  \begin{block}{Problématique}
    Les terminaux constituent des points d'entrée sensibles, tandis que les antivirus
    basés sur les signatures restent insuffisants face aux variantes nouvelles ou inconnues.
    La nécessité de détecter automatiquement un comportement ransomware sur un parc de
    machines hétérogène --- et de déclencher une réponse maîtrisée
    \textbf{sans déployer de vrai malware} --- constitue le défi central de ce travail.
  \end{block}

\end{frame}
\note{
  \textbf{[CONTEXTE \& PROBLÉMATIQUE — 1:00]}\\[4pt]
  Les organisations dépendent aujourd'hui fortement de leurs systèmes informatiques,
  ce qui les expose de plus en plus aux cyberattaques — en particulier au ransomware,
  ce logiciel malveillant qui chiffre les données puis exige une rançon.
  Selon l'ENISA, ces attaques ont augmenté de 73\,\% en 2023~; sans outil dédié,
  il faut en moyenne 21 jours pour les détecter, pour un coût qui dépasse souvent
  le million de dollars. Le problème, c'est que les antivirus classiques, basés
  sur des signatures, ne suffisent plus face à des comportements inconnus.
  D'où la question centrale de ce travail~: comment détecter et répondre à un
  comportement ransomware, dans un environnement contrôlé, sans déployer de vrai malware~?
}

% ─── Objectifs ────────────────────────────────────────────────────────────
\begin{frame}{Introduction Générale}

  \begin{block}{Objectif général}
    Concevoir et mettre en œuvre un \textbf{système de détection et de réponse
    aux attaques ransomware} basé sur la surveillance comportementale des fichiers
    et des processus, avec une console de supervision traçable.
  \end{block}

  \vspace{4pt}

  \begin{block}{Objectifs spécifiques}
    \begin{itemize}\setlength\itemsep{2pt}
      \item Surveiller les événements fichiers en temps réel sur les terminaux
      \item Analyser les comportements et calculer un score de risque configurable
      \item Générer automatiquement des alertes et des incidents de sécurité
      \item Proposer et exécuter des actions de protection avec \textbf{contrôle humain}
      \item Sécuriser la console par une authentification forte (2FA TOTP RFC\,6238)
      \item Valider le pipeline à l'aide de scénarios d'attaque représentatifs
    \end{itemize}
  \end{block}

\end{frame}
\note{
  \textbf{[OBJECTIFS — 2:30]}\\[4pt]
  L'objectif général est de concevoir un système de détection et de réponse aux
  attaques ransomware, basé sur la surveillance comportementale des fichiers et
  des processus. Concrètement, il s'agit de surveiller les événements en temps réel,
  de calculer un score de risque, de générer automatiquement des alertes et des
  incidents, de proposer des actions de protection avec un contrôle humain
  systématique, de sécuriser la console par une double authentification, et de
  valider le tout sur des scénarios d'attaque réalistes.
}

% ─── Organisation ─────────────────────────────────────────────────────────
\begin{frame}{Organisation du mémoire}

  \vspace{0.2cm}

  \begin{columns}[T]
    \column{0.32\textwidth}
      \begin{block}{Chapitre I\\Revue de littérature}
        \small Sécurité informatique, logiciels malveillants, ransomwares,
        antivirus, IDS, EDR~: positionnement de RansomShield.
      \end{block}
    \column{0.32\textwidth}
      \begin{block}{Chapitre II\\Analyse et Conception}
        \small Besoins fonctionnels et non fonctionnels, architecture générale,
        modèle de données, processus de détection et mesures de sécurité.
      \end{block}
    \column{0.32\textwidth}
      \begin{block}{Chapitre III\\Réalisation et Résultats}
        \small Implémentation des composants, scénarios de validation,
        résultats obtenus, limites et perspectives.
      \end{block}
  \end{columns}

  \vspace{8pt}

  \begin{tcolorbox}[
    colback=light,colframe=primary,arc=5pt,boxrule=0.8pt,halign=center,
    left=8pt,right=8pt,top=5pt,bottom=5pt
  ]
    \small\color{textdark}
    Mémoire rédigé en \LaTeX{} selon le canevas IFRI ---
    Licence Informatique, option Sécurité Informatique --- Année 2023--2024
  \end{tcolorbox}

\end{frame}
\note{
  \textbf{[ORGANISATION DU MÉMOIRE — 3:15]}\\[4pt]
  Le mémoire suit le canevas IFRI en trois chapitres~: la revue de littérature,
  l'analyse et la conception du système, puis sa réalisation et ses résultats.
  Voyons d'abord les fondements théoriques.
}

% =========================================================================
\section{Revue de littérature}
% =========================================================================

\sectionslide{02}{Revue de littérature}{book}

% ─── Définition + comportements ───────────────────────────────────────────
\begin{frame}{Revue de littérature}

  \begin{block}{Définition d'un ransomware}
    Un ransomware est un logiciel malveillant qui vise à bloquer l'accès à un
    système ou à des données, souvent par chiffrement, puis à réclamer une rançon
    à la victime. \textit{(NIST SP\,1800-25, CISA)}
  \end{block}

  \vspace{4pt}

  \begin{block}{Comportements caractéristiques --- détectables avant les dommages irréversibles}
    \begin{itemize}\setlength\itemsep{2pt}
      \item Renommage massif de fichiers avec extension suspecte
            (\texttt{.locked}, \texttt{.enc}, \texttt{.crypt}...)
      \item Création d'une note de rançon (\texttt{HOW\_TO\_DECRYPT.txt})
      \item Suppression des clichés instantanés VSS
            (\texttt{vssadmin delete shadows /all})
      \item Utilisation d'outils légitimes détournés ---
            LOLBins (\texttt{bcdedit}, \texttt{wmic}, \texttt{vssadmin})
      \item Suppression massive de fichiers, accès réseau inhabituel
    \end{itemize}
    Ces signaux constituent la base de la \textbf{détection comportementale}.
  \end{block}

\end{frame}
\note{
  \textbf{[DÉFINITION \& COMPORTEMENTS — 3:50]}\\[4pt]
  Selon le NIST, un ransomware est un logiciel malveillant qui bloque l'accès aux
  données, le plus souvent par chiffrement, pour ensuite réclamer une rançon.
  Ce qui compte pour nous, c'est que ce type d'attaque laisse des signes avant que
  les dégâts ne soient irréversibles~: renommage massif de fichiers, apparition
  d'extensions suspectes comme \texttt{.locked}, création d'une note de rançon,
  suppression des sauvegardes système, ou encore utilisation détournée d'outils
  légitimes. Ce sont ces signaux qui constituent la base de notre détection
  comportementale.
}

% ─── Solutions existantes ─────────────────────────────────────────────────
\begin{frame}{Revue de littérature}

  \begin{block}{Quelques solutions existantes}
    \begin{itemize}\setlength\itemsep{2pt}
      \item \textbf{Antivirus} --- détection par signatures statiques (ClamAV, Windows Defender)
      \item \textbf{SIEM génériques} --- collecte de journaux sans corrélation comportementale
      \item \textbf{EDR commerciaux} --- CrowdStrike Falcon, SentinelOne, Microsoft Defender XDR
    \end{itemize}
  \end{block}

  \vspace{4pt}

  \begin{block}{Limites des solutions existantes}
    \begin{itemize}\setlength\itemsep{2pt}
      \item \bad Contournables par variantes nouvelles, obfuscation ou LOLBins
      \item \bad Coût élevé des EDR commerciaux (hors portée académique)
      \item \bad Absence d'outil open-source académique démontrable en conditions réelles
    \end{itemize}
  \end{block}

  \vspace{4pt}

  \begin{alertblock}{Notre réponse : RansomShield}
    Approche comportementale open-source, multi-plateforme, démontrable sur
    machines virtuelles réelles --- \textbf{contrôle humain à chaque étape}.
  \end{alertblock}

\end{frame}
\note{
  \textbf{[SOLUTIONS EXISTANTES — 5:15]}\\[4pt]
  Les solutions existantes ont chacune leurs limites. Les antivirus classiques
  détectent par signatures, donc ils passent à côté des nouvelles variantes.
  Les SIEM collectent des journaux mais sans corrélation comportementale spécifique.
  Les EDR commerciaux, comme CrowdStrike ou SentinelOne, sont efficaces mais hors
  de portée pour un contexte académique. C'est dans cet espace que se positionne
  RansomShield~: une approche comportementale, open-source, avec un contrôle
  humain garanti à chaque étape. Voyons maintenant comment ce système a été conçu.
}

% =========================================================================
\section{Analyse et Conception}
% =========================================================================

\sectionslide{03}{Analyse et Conception}{drafting-compass}

% ─── Architecture ─────────────────────────────────────────────────────────
\begin{frame}{Architecture générale de RansomShield}
  \centering
  \vspace{0.3cm}
  \begin{adjustbox}{max width=\textwidth,max totalheight=0.72\textheight,center}
  \begin{tikzpicture}[
    box/.style={
      rectangle, rounded corners=6pt, draw=primary, line width=1.2pt,
      fill=light, minimum width=3.9cm, minimum height=2.6cm,
      align=center, font=\normalsize
    }
  ]
    \node[box] (b1) at (0,0)
      {\Large\faDesktop\\[4pt]\textbf{Agents}\\ \small Windows / Linux / macOS};
    \node[box] (b2) at (5.2,0)
      {\Large\faServer\\[4pt]\textbf{API Laravel}\\ \small Clé par agent};
    \node[box] (b3) at (10.4,0)
      {\Large\faBrain\\[4pt]\textbf{Moteur de détection}\\ \small Score de risque};
    \node[box] (b4) at (15.6,0)
      {\Large\faChartBar\\[4pt]\textbf{Console SOC}\\ \small Alertes \& incidents};

    \draw[-{Stealth[length=3.5mm]}, primary, line width=1.4pt] (b1) -- (b2);
    \draw[-{Stealth[length=3.5mm]}, primary, line width=1.4pt] (b2) -- (b3);
    \draw[-{Stealth[length=3.5mm]}, primary, line width=1.4pt] (b3) -- (b4);

    \node[box, fill=primary!10, draw=primary!70!black, minimum width=6.5cm, minimum height=1.6cm]
      (op) at (7.8,-3.6)
      {\Large\faUserCheck\quad\textbf{Opérateur humain}\\ \small approuve chaque action avant exécution};

    \draw[-{Stealth[length=3.5mm]}, primary!70!black, dashed, line width=1.2pt, rounded corners=8pt]
      (15.6,-1.3) -- (15.6,-2.4) -- (9.8,-2.4) -- (9.8,-2.8);
    \draw[-{Stealth[length=3.5mm]}, primary!70!black, dashed, line width=1.2pt, rounded corners=8pt]
      (5.8,-2.8) -- (5.8,-2.4) -- (0,-2.4) -- (0,-1.3);
  \end{tikzpicture}
  \end{adjustbox}
  \vspace{6pt}
  {\small\color{muted}
   Aucune action n'est exécutée automatiquement — l'opérateur valide toujours avant.}
\end{frame}
\note{
  \textbf{[ARCHITECTURE — 6:15]}\\[4pt]
  L'architecture repose sur un flux simple. Les agents Python, installés sur les
  machines surveillées, envoient les événements à l'API Laravel via une clé propre
  à chaque agent. Le moteur de détection analyse ces événements, calcule un score
  de risque, et déclenche une alerte puis un incident si le seuil est dépassé.
  L'opérateur consulte la console SOC et valide chaque action avant son exécution~:
  c'est notre règle d'or, aucune action automatique sans validation humaine.
}

% ─── Les 4 composants ─────────────────────────────────────────────────────
\begin{frame}{Les quatre composants du système}

  \begin{columns}[T]
    \column{0.48\textwidth}
      \begin{block}{\faCode\quad Agent Python}
        \small Surveille les dossiers via \textbf{Watchdog / inotify}.
        File locale SQLite (résilience réseau).
        Enrôlement par token à usage unique.
        Service \texttt{systemd} --- démarrage automatique.
      \end{block}
      \vspace{4pt}
      \begin{block}{\faNetworkWired\quad Découverte réseau}
        \small fping (ICMP) + ARP sweep sur les \texttt{/24}.
        Scan automatique toutes les 5 minutes.
        Aucun enrôlement sans \textbf{validation manuelle admin}.
      \end{block}

    \column{0.48\textwidth}
      \begin{block}{\faDesktop\quad Console SOC Laravel}
        \small Dashboard temps réel (Chart.js, AJAX 30\,s).
        Alertes, incidents, timeline, file d'approbation.
        Notifications~: web, e-mail, son, webhook.
      \end{block}
      \vspace{4pt}
      \begin{block}{\faFlask\quad Module de simulation}
        \small 5 scénarios d'attaque, jusqu'à \textbf{22 événements}.
        Tous marqués \texttt{is\_simulation=true} ---
        nettoyage facile après chaque démo.
      \end{block}
  \end{columns}

\end{frame}
\note{
  \textbf{[LES 4 COMPOSANTS — 7:15]}\\[4pt]
  Le système repose sur quatre composants. L'agent Python surveille les fichiers
  en temps réel et conserve une file locale pour ne rien perdre en cas de coupure
  réseau. La console SOC centralise le suivi~: alertes, incidents, et validation
  des actions. La découverte réseau scanne automatiquement les sous-réseaux, mais
  aucun agent ne peut s'enrôler sans validation manuelle de l'administrateur.
  Enfin, le module de simulation permet de rejouer des scénarios d'attaque
  réalistes, jusqu'à vingt-deux événements, sans jamais toucher à un vrai malware.
}

% ─── Diagramme de cas d'utilisation (UML TikZ) ───────────────────────────
\begin{frame}{Diagramme de cas d'utilisation}
\vspace{-0.15cm}
\begin{center}
\begin{adjustbox}{max width=0.9\textwidth,max totalheight=0.78\textheight,center}
\begin{tikzpicture}[
  uc/.style={
    ellipse, draw=primary, fill=light,
    font=\large\bfseries, align=center,
    text width=5.0cm, minimum height=1.1cm,
    inner sep=4pt, text=textdark
  },
  ucshared/.style={
    ellipse, draw=primary!70!black, fill=primary!15,
    font=\large\bfseries, align=center,
    text width=5.0cm, minimum height=1.1cm,
    inner sep=4pt, text=textdark
  }
]

%% ── ACTEUR 1 : Administrateur (gauche haut) ──
\draw[fill=darkbg, draw=darkbg, line width=1.4pt]
  (-0.7, 2.85) circle (0.38);
\draw[darkbg, line width=1.5pt]
  (-0.7, 2.47)--(-0.7, 1.65)
  (-1.25, 2.15)--(-0.1, 2.15)
  (-0.7, 1.65)--(-1.2, 0.95)
  (-0.7, 1.65)--(-0.2, 0.95);
\node[below, font=\large\bfseries, text=darkbg, align=center]
  at (-0.7, 0.75) {Administrateur};
\coordinate (adminPt) at (-0.05, 2.15);

%% ── ACTEUR 2 : Analyste SOC (gauche bas) ──
\draw[fill=primary!75!black, draw=primary!75!black, line width=1.4pt]
  (-0.7, -2.35) circle (0.38);
\draw[primary!75!black, line width=1.5pt]
  (-0.7,-2.73)--(-0.7,-3.55)
  (-1.25,-3.1)--(-0.1,-3.1)
  (-0.7,-3.55)--(-1.2,-4.25)
  (-0.7,-3.55)--(-0.2,-4.25);
\node[below, font=\large\bfseries, text=primary!75!black, align=center]
  at (-0.7,-4.45) {Analyste SOC};
\coordinate (socPt) at (-0.05, -3.1);

%% ── FRONTIÈRE SYSTÈME ──
\draw[draw=primary, line width=1.8pt, rounded corners=8pt, fill=primary!3]
  (1.5, 4.35) rectangle (11.8, -4.7);
\node[font=\large\bfseries, text=primary, fill=white, inner sep=6pt,
      rounded corners=5pt, draw=primary, line width=1pt]
  at (6.65, 4.35) {$\ll$\,system\,$\gg$\quad\faIcon{shield-alt}\enspace RansomShield};

%% ── CAS D'UTILISATION (6, colonne unique) ──
\node[uc]       (uc1) at (6.5,  3.25)  {Gérer agents \& réseaux};
\node[uc]       (uc2) at (6.5,  1.95)  {Configurer le système};
\node[uc]       (uc3) at (6.5,  0.65)  {Lancer une simulation};
\node[ucshared] (uc4) at (6.5, -0.65)  {Visualiser alertes \& incidents};
\node[ucshared] (uc5) at (6.5, -1.95)  {Approuver / rejeter\\une action};
\node[ucshared] (uc6) at (6.5, -3.25)  {Consulter la timeline\\d'audit};

%% ── ASSOCIATIONS ADMIN (colonne unique, ligne directe, aucun obstacle) ──
\foreach \uc in {uc1,uc2,uc3,uc4,uc5,uc6}
  \draw[gray!55, thin] (adminPt) -- (\uc.west);

%% ── ASSOCIATIONS SOC (uniquement les cas partagés) ──
\foreach \uc in {uc4,uc5,uc6}
  \draw[primary!55, thin, dashed] (socPt) -- (\uc.west);

%% ── LÉGENDE ──
\draw[draw=muted!60, rounded corners=4pt, fill=white]
  (2.5, -4.15) rectangle (11.0, -4.5);
\node[font=\small, text=textdark, align=center] at (6.75, -4.325) {
  \tikz{\draw[primary,fill=light](0,0)ellipse(0.28 and 0.15);}\ \ Admin seul
  \qquad
  \tikz{\draw[primary!70!black,fill=primary!15](0,0)ellipse(0.28 and 0.15);}\ \ Partagé SOC
};

\end{tikzpicture}
\end{adjustbox}
\end{center}
\end{frame}
\note{
  \textbf{[DIAGRAMME UML — 8:45]}\\[4pt]
  Le système distingue deux acteurs. L'administrateur gère les agents, configure
  le système et lance les simulations. L'analyste SOC, lui, se concentre sur le
  traitement opérationnel~: il visualise les alertes et les incidents, approuve
  ou rejette les actions de protection, et consulte la timeline d'audit — mais
  n'a pas accès à la configuration.
}

% ─── Modèle de données ────────────────────────────────────────────────────
\begin{frame}{Modèle de données}

  \begin{columns}[T]
    \column{0.48\textwidth}
      \begin{block}{\faDatabase\quad Tables principales}
        \small
        \renewcommand{\arraystretch}{1.2}
        \begin{tabular}{@{}l l@{}}
          \texttt{agents}              & enrôlement, état, risque \\
          \texttt{events}              & événements fichiers       \\
          \texttt{alerts}              & score, niveau, signaux    \\
          \texttt{incidents}           & fiche, statut, timeline   \\
          \texttt{detection\_rules}    & règles configurables      \\
          \texttt{protection\_actions} & décisions SOC             \\
          \texttt{audit\_logs}         & traçabilité horodatée     \\
          \texttt{simulation\_runs}    & profils d'attaque         \\
        \end{tabular}
      \end{block}

    \column{0.48\textwidth}
      \begin{block}{\faProjectDiagram\quad Relations clés}
        \small
        \begin{itemize}\setlength\itemsep{3pt}
          \item \texttt{agents} $\to$ \texttt{events} (1--N)
          \item \texttt{events} $\to$ \texttt{alerts} (N--1)
          \item \texttt{alerts} $\to$ \texttt{incidents} (N--1)
          \item \texttt{incidents} $\to$ \texttt{protection\_actions} (1--N)
          \item \texttt{incidents} $\to$ \texttt{audit\_logs} (1--N)
          \item \texttt{simulation\_runs} $\to$ \texttt{events} (1--N)
        \end{itemize}
      \end{block}
      \vspace{4pt}
      \begin{tcolorbox}[
        colback=light,colframe=primary,arc=4pt,boxrule=0.6pt,
        left=6pt,right=6pt,top=4pt,bottom=4pt
      ]
        \tiny\color{textdark}
        8 tables MySQL --- pipeline complet~:
        événement $\to$ alerte $\to$ incident $\to$ action $\to$ audit
      \end{tcolorbox}
  \end{columns}

\end{frame}
\note{
  \textbf{[MODÈLE DE DONNÉES — 9:45]}\\[4pt]
  Le modèle de données repose sur huit tables MySQL. Le pipeline se lit de gauche
  à droite~: un agent génère des événements, qui déclenchent des alertes, qui
  créent des incidents, qui donnent lieu à des actions de protection — le tout
  tracé dans un journal d'audit. Chaque action est horodatée avec l'identité de
  l'opérateur qui l'a validée. Passons maintenant à la réalisation concrète.
}

% =========================================================================
\section{Réalisation et Résultats}
% =========================================================================

\sectionslide{04}{Réalisation et Résultats}{cogs}

% ─── Choix techniques ─────────────────────────────────────────────────────
\begin{frame}{Choix techniques}

  \begin{columns}[T]
    \column{0.48\textwidth}
      \begin{block}{Backend \& Base de données}
        \small
        \begin{tabular}{@{}l@{\enspace}l@{}}
          \faServer   & \textbf{Laravel 11} / PHP 8.3 \\[3pt]
          \faDatabase & \textbf{MySQL 8.0} (Eloquent ORM) \\[3pt]
          \faCode     & API REST sécurisée (clé per-agent) \\
        \end{tabular}
      \end{block}
      \vspace{4pt}
      \begin{block}{Agent de surveillance}
        \small
        \begin{tabular}{@{}l@{\enspace}l@{}}
          \faCode      & \textbf{Python 3.10} \\[3pt]
          \faSearch    & \textbf{Watchdog} (inotify Linux) \\[3pt]
          \faDatabase  & \textbf{SQLite} (file locale offline) \\[3pt]
          \faMicrochip & psutil (processus) \\
        \end{tabular}
      \end{block}

    \column{0.48\textwidth}
      \begin{block}{Frontend}
        \small
        \begin{tabular}{@{}l@{\enspace}l@{}}
          \faDesktop  & \textbf{Blade} (templates Laravel) \\[3pt]
          \faChartBar & \textbf{Chart.js} (intégré local) \\[3pt]
          \faSync     & AJAX toutes les 30\,s \\
        \end{tabular}
      \end{block}
      \vspace{4pt}
      \begin{block}{Infrastructure de test}
        \small
        \begin{tabular}{@{}l@{\enspace}l@{}}
          \faServer        & 3 VMs \textbf{KVM/libvirt} \\[3pt]
          \faNetworkWired  & Réseau \texttt{10.20.0.0/24} \\[3pt]
          \faLock          & \textbf{2FA TOTP} RFC\,6238 \\[3pt]
          \faKey           & Token usage unique (enrôlement) \\
        \end{tabular}
      \end{block}
  \end{columns}

\end{frame}
\note{
  \textbf{[CHOIX TECHNIQUES — 11:00]}\\[4pt]
  Côté technique, le backend repose sur Laravel 11 et PHP 8.3, avec une base MySQL
  et une API sécurisée par une clé propre à chaque agent. L'agent de surveillance
  est écrit en Python avec la bibliothèque Watchdog, et conserve ses événements
  dans une file SQLite locale. Le frontend utilise Chart.js en local, sans
  dépendance externe. Les tests ont été menés sur trois machines virtuelles KVM,
  avec une authentification à double facteur.
}

% ─── Agent Python + Moteur ────────────────────────────────────────────────
\begin{frame}{Agent Python et moteur de détection}

  \begin{columns}[T]
    \column{0.44\textwidth}
      \begin{block}{Cycle de vie de l'agent}
        \small
        \begin{enumerate}\setlength\itemsep{2pt}
          \item Démarrage : lecture \texttt{.env}
          \item Enrôlement : token $\to$ clé API
          \item Watchdog observe les dossiers
          \item Événement $\to$ file SQLite locale
          \item \texttt{POST /api/agent/events}
          \item Heartbeat toutes les \textbf{30\,s}
        \end{enumerate}
        \vspace{5pt}
        \textit{\tiny Installation en une commande (Linux, Windows, macOS)~:}\\
        {\color{textdark!60}\tiny\ttfamily curl -sSL http://<soc>/install.sh | bash}
      \end{block}

    \column{0.52\textwidth}
      \begin{block}{Scoring comportemental}
        \small
        \renewcommand{\arraystretch}{1.2}
        \begin{tabular}{@{}l r@{}}
          \toprule
          \textbf{Signal détecté} & \textbf{Score} \\
          \midrule
          Extension suspecte (\texttt{.locked}) & $+80$ \\
          Note de rançon créée                  & $+55$ \\
          Suppression clichés VSS               & $+60$ \\
          LOLBin suspect                        & $+45$ \\
          Suppression massive fichiers          & $+50$ \\
          \bottomrule
        \end{tabular}
        \vspace{5pt}

        \begin{tabular}{@{}l c@{}}
          Score $<$ 30  & \textcolor{norm}{\textbf{Normal}}   \\
          30 à 59       & \textcolor{high}{\textbf{Suspect}}  \\
          60 à 99       & \textcolor{high}{\textbf{Élevé}}   \\
          $\geq$ 100    & \textcolor{crit}{\textbf{Critique}} \\
        \end{tabular}
        \vspace{3pt}
        {\tiny\color{muted} Alerte + incident créés automatiquement en $<$\,5\,s}
      \end{block}
  \end{columns}

\end{frame}
\note{
  \textbf{[AGENT \& MOTEUR — 12:15]}\\[4pt]
  Au démarrage, l'agent s'enrôle via un token à usage unique qui lui donne une clé
  API permanente, puis il surveille les dossiers avec Watchdog et envoie un signal
  de présence toutes les trente secondes. Le moteur de détection calcule un score
  cumulatif configurable. Par exemple~: un fichier renommé en \texttt{.locked}
  vaut 80 points, une note de rançon 55 points~; le total de 135 déclenche
  immédiatement une alerte critique et la création d'un incident. Tout ce
  processus prend moins de cinq secondes.
}

% ─── Console SOC ──────────────────────────────────────────────────────────
\begin{frame}{Console SOC --- Interface de supervision}

  \begin{columns}[T]
    \column{0.40\textwidth}
      \begin{block}{Pages de la console}
        \small
        \begin{tabular}{@{}l@{\enspace}l@{}}
          \faHome        & Tableau de bord \\
          \faRobot       & Agents (état, heartbeat) \\
          \faBell        & Alertes (score, signaux) \\
          \faFolderOpen  & Incidents \\
          \faHistory     & Timeline d'audit \\
          \faCheckSquare & File d'approbation \\
          \faGavel       & Règles (CRUD live) \\
          \faSearch      & Recherche globale \\
          \faBell        & Notifications \\
          \faCog         & Paramètres système \\
        \end{tabular}
      \end{block}

    \column{0.57\textwidth}
      \begin{tcolorbox}[
        colback=darkbg,colframe=primary,arc=5pt,
        boxrule=0.8pt,left=2pt,right=2pt,top=2pt,bottom=2pt
      ]
        \includegraphics[width=\linewidth]{dashboard}
      \end{tcolorbox}
      \vspace{2pt}
      {\tiny\color{muted}%
        \faSync\ Mise à jour AJAX toutes les 30\,s \quad
        \faBell\ Web · E-mail SMTP · Son · Webhook}
  \end{columns}

\end{frame}
\note{
  \textbf{[CONSOLE SOC — 13:30]}\\[4pt]
  La console SOC centralise dix pages fonctionnelles, avec un tableau de bord mis
  à jour automatiquement toutes les trente secondes. Le point important, c'est que
  toute action de protection passe par une file d'approbation~: l'analyste doit
  explicitement valider ou rejeter chaque action avant qu'elle ne s'exécute.
  Les notifications partent simultanément par quatre canaux~: le web, l'e-mail,
  le son, et un webhook. Vous verrez tout cela dans la démonstration.
}

% ─── Tests de validation ──────────────────────────────────────────────────
\begin{frame}{Tests de validation et résultats}

  \begin{columns}[T]
    \column{0.50\textwidth}
      \begin{block}{5 scénarios d'attaque testés}
        \small
        \renewcommand{\arraystretch}{1.2}
        \begin{tabular}{@{}l c c@{}}
          \toprule
          \textbf{Scénario} & \textbf{Évts} & \textbf{Résultat} \\
          \midrule
          Chiffrement basique  &  7 & \ok \\
          Kill chain complète  & 22 & \ok \\
          Chiffrement massif   & 18 & \ok \\
          Exfiltration données &  8 & \ok \\
          LOLBins suspects     &  5 & \ok \\
          \bottomrule
        \end{tabular}
        \vspace{3pt}
        {\tiny\color{muted} Environnement~: 3 VMs KVM --- réseau \texttt{10.20.0.0/24}}
      \end{block}

    \column{0.47\textwidth}
      \begin{block}{Résultats obtenus}
        \small
        \begin{itemize}\setlength\itemsep{2pt}
          \item[\ok] Détection en \textbf{$<$ 5 secondes} (tous scénarios)
          \item[\ok] Aucun événement perdu (file SQLite)
          \item[\ok] Pipeline complet validé :
                     événement $\to$ alerte $\to$ incident $\to$ timeline
          \item[\ok] Rollback isolation réseau fonctionnel
          \item[\ok] 2FA TOTP et API sécurisée (per-agent)
        \end{itemize}
      \end{block}
      \vspace{3pt}
      \begin{alertblock}{Limite identifiée}
        \small Tests conduits sur Linux uniquement.
        Un attaquant ralentissant son rythme réduit la visibilité.
      \end{alertblock}
  \end{columns}

\end{frame}
\note{
  \textbf{[TESTS \& RÉSULTATS — 14:45]}\\[4pt]
  Nous avons validé le système sur cinq scénarios d'attaque, de sept à vingt-deux
  événements. Tous ont été détectés en moins de cinq secondes, sans perte
  d'événement, avec un pipeline complet validé de bout en bout, de l'événement
  jusqu'à la timeline. Je précise honnêtement une limite~: les tests ont été
  conduits uniquement sur Linux, et un attaquant qui ralentirait son rythme
  resterait moins visible. Passons maintenant à la démonstration du système.
}

% ─────────────────────────────────────────────────────────────────────────
%  SLIDE TRANSITION DÉMONSTRATION
% ─────────────────────────────────────────────────────────────────────────
{
\setbeamertemplate{headline}{}
\setbeamertemplate{footline}{}
\begin{frame}[plain,noframenumbering]
  \begin{tikzpicture}[overlay,remember picture]
    \fill[primary]
      (current page.north west) rectangle (current page.south east);
    \node[anchor=center,font=\Huge\bfseries\color{white}]
      at ([yshift=1cm]current page.center)
      {\faPlay\quad Démonstration Live};
    \node[anchor=center,font=\normalsize\color{white!80}]
      at ([yshift=-0.5cm]current page.center)
      {5 minutes --- système réel sur VMs KVM};
  \end{tikzpicture}
\end{frame}
}
\note{
  \textbf{[DÉMO LIVE — 15:00, checklist 5 min, hors chrono présentation]}\\[4pt]
  1. Ouvrir \texttt{http://10.20.0.1} (console SOC)\\[3pt]
  2. Se connecter : identifiants + \textbf{code 2FA}\\[3pt]
  3. Vérifier : 3 agents en ligne (heartbeat vert)\\[3pt]
  4. \textbf{Agents} $\to$ VM-02 $\to$ \textbf{Lancer simulation} $\to$ "Kill chain complète" (22 évts)\\[3pt]
  5. \textbf{Dashboard} : le score monte en temps réel\\[3pt]
  6. \textbf{Alertes} : niveau Critique + signaux détectés\\[3pt]
  7. \textbf{Incidents} : ouvrir l'incident auto-créé\\[3pt]
  8. \textbf{File d'approbation} : approuver l'isolation réseau\\[3pt]
  9. \textbf{Timeline} : tout est horodaté et tracé\\[3pt]
  \textit{Rester calme, ne pas se précipiter — chaque étape a le temps qu'il faut.}
}

% =========================================================================
\section{Conclusion et Perspectives}
% =========================================================================

\sectionslide{05}{Conclusion et Perspectives}{rocket}

% ─── Conclusion ───────────────────────────────────────────────────────────
\begin{frame}{Conclusion et Perspectives}

  \begin{block}{Conclusion}
    \small RansomShield est une solution opérationnelle couvrant l'ensemble du cycle
    de traitement d'un incident~:
    \textbf{collecte $\to$ analyse $\to$ alerte $\to$ réponse $\to$ audit}.
    \begin{itemize}\setlength\itemsep{1pt}
      \item[\ok] Agent Python multi-plateforme --- installation en une commande
      \item[\ok] Console SOC~: simulation, 4 canaux de notification, rapports PDF, 2FA
      \item[\ok] Moteur de détection~: scoring configurable, règles CRUD, $<$\,5\,s
      \item[\ok] Environnement contrôlé~: aucun vrai malware déployé
    \end{itemize}
  \end{block}

  \vspace{4pt}

  \begin{block}{Perspectives d'évolution}
    \small
    \begin{itemize}\setlength\itemsep{1pt}
      \item[\arr] \textbf{HTTPS} --- chiffrement TLS en production
      \item[\arr] \textbf{RBAC} --- rôles distincts (administrateur, analyste SOC)
      \item[\arr] \textbf{Machine Learning} --- Isolation Forest (baseline comportementale)
      \item[\arr] \textbf{MITRE ATT\&CK} --- alignement et couverture du référentiel
      \item[\arr] \textbf{Intégration SIEM} --- export Syslog vers Splunk / ELK Stack
    \end{itemize}
  \end{block}

\end{frame}
\note{
  \textbf{[CONCLUSION \& PERSPECTIVES — 20:00]}\\[4pt]
  En résumé, RansomShield est une solution fonctionnelle qui couvre tout le cycle
  de traitement d'un incident~: collecte, analyse, alerte, réponse et audit.
  Elle repose sur un agent multi-plateforme, une console complète, un moteur de
  détection réactif, et un contrôle humain à chaque étape. Ce travail a ses
  limites~: il a été testé dans un environnement contrôlé, avec des réponses en
  partie simulées. Les perspectives sont claires~: ajouter le chiffrement HTTPS,
  des rôles distincts pour les utilisateurs, puis à terme du machine learning
  pour affiner la détection, et une intégration avec les outils SIEM du marché.
  Je vous remercie pour votre attention et reste à votre disposition pour vos
  questions.
}

% ─────────────────────────────────────────────────────────────────────────
%  SLIDE FINAL — MERCI / QUESTIONS
% ─────────────────────────────────────────────────────────────────────────
{
\setbeamertemplate{headline}{}
\setbeamertemplate{footline}{}
\begin{frame}[plain,noframenumbering]

  \begin{tikzpicture}[overlay,remember picture]
    \shade[top color=darkbg,bottom color=primary!80!darkbg]
      (current page.north west) rectangle (current page.south east);
  \end{tikzpicture}

  \centering\vspace{0.8cm}

  {\color{white}\LARGE\bfseries Merci de votre aimable attention\,!}

  \vspace{0.6cm}

  {\color{primary!40!white}\fontsize{52}{52}\selectfont ?}

  \vspace{0.6cm}

  \begin{tcolorbox}[
    colback=white!8!darkbg,colframe=primary,
    arc=8pt,boxrule=1.2pt,width=0.70\textwidth,center,
    halign=center,left=12pt,right=12pt,top=10pt,bottom=10pt
  ]
    \color{white}\small
    \textbf{Codjo Fréjus Loïc SANGRONIO}\\[4pt]
    \texttt{lsangronio@gmail.com} \textbar{} IFRI-UAC 2023--2024\\[5pt]
    {\scriptsize Encadreur~: \textbf{Ing. Guy KPEDJO}}
  \end{tcolorbox}

\end{frame}
}
\note{
  \textbf{[QUESTIONS DU JURY — aide-mémoire]}\\[4pt]
  \textbullet\ \ \textbf{Pourquoi Python et pas C/Go ?} $\to$ Watchdog exploite inotify nativement, CPU minimal, portable Linux/Windows/macOS\\[3pt]
  \textbullet\ \ \textbf{Faux positifs ?} $\to$ seuils configurables, \textbf{opérateur humain} valide ou rejette chaque action\\[3pt]
  \textbullet\ \ \textbf{Pourquoi pas HTTPS ?} $\to$ LAN contrôlé ; c'est la 1\textsuperscript{re} perspective pour la prod\\[3pt]
  \textbullet\ \ \textbf{Différence avec un antivirus ?} $\to$ comportement, pas signature ; fonctionne sur variantes inconnues\\[3pt]
  \textbullet\ \ \textbf{Testé sur Windows ?} $\to$ Non, limite assumée : Linux uniquement\\[3pt]
  \textit{Garder les \textbf{annexes} prêtes (file d'approbation, captures d'écran).}
}

% =========================================================================
%  ANNEXES (hors numérotation principale)
% =========================================================================
\appendix

\begin{frame}{Annexe A --- File d'approbation et rollback}
  \begin{columns}[T]
    \column{0.44\textwidth}
      \begin{block}{Principe}
        \small Le moteur \textbf{propose} --- l'opérateur \textbf{décide}.\\[4pt]
        Chaque action affiche~:
        \begin{itemize}\setlength\itemsep{1pt}
          \item Type (scan, isolation, quarantaine)
          \item Niveau de risque de l'incident
          \item Agent concerné
        \end{itemize}
        \vspace{3pt}
        Boutons \textbf{Approuver / Rejeter} --- tout est horodaté.\\[3pt]
        \textbf{Rollback} possible sur l'isolation réseau.
      \end{block}
    \column{0.52\textwidth}
      \begin{tcolorbox}[colback=darkbg,colframe=primary,arc=3pt,
          boxrule=0.8pt,left=2pt,right=2pt,top=2pt,bottom=2pt]
        \includegraphics[width=\linewidth]{protection_actions}
      \end{tcolorbox}
  \end{columns}
\end{frame}

\begin{frame}{Annexe B --- Captures d'écran}
  \begin{columns}[T]
    \column{0.48\textwidth}
      {\scriptsize\bfseries\color{primary}\faBell\quad Alertes}
      \vspace{2pt}
      \begin{tcolorbox}[colback=darkbg,colframe=primary,arc=2pt,
          boxrule=0.6pt,left=1pt,right=1pt,top=1pt,bottom=1pt]
        \includegraphics[width=\linewidth]{alerts}
      \end{tcolorbox}
      \vspace{3pt}
      {\scriptsize\bfseries\color{primary}\faHistory\quad Timeline incident}
      \vspace{2pt}
      \begin{tcolorbox}[colback=darkbg,colframe=primary,arc=2pt,
          boxrule=0.6pt,left=1pt,right=1pt,top=1pt,bottom=1pt]
        \includegraphics[width=\linewidth]{incident_timeline}
      \end{tcolorbox}
    \column{0.48\textwidth}
      {\scriptsize\bfseries\color{primary}\faRobot\quad Agents}
      \vspace{2pt}
      \begin{tcolorbox}[colback=darkbg,colframe=primary,arc=2pt,
          boxrule=0.6pt,left=1pt,right=1pt,top=1pt,bottom=1pt]
        \includegraphics[width=\linewidth]{agents}
      \end{tcolorbox}
      \vspace{3pt}
      {\scriptsize\bfseries\color{primary}\faFlask\quad Simulation en cours}
      \vspace{2pt}
      \begin{tcolorbox}[colback=darkbg,colframe=primary,arc=2pt,
          boxrule=0.6pt,left=1pt,right=1pt,top=1pt,bottom=1pt]
        \includegraphics[width=\linewidth]{simulation}
      \end{tcolorbox}
  \end{columns}
\end{frame}

\end{document}
